% part: intuitionistic-logic
% chapter: propositions-as-types
% section: normalization

\documentclass[../../../include/open-logic-section]{subfiles}

\begin{document}

\olfileid{int}{pty}{nor}

\olsection{نرمال‌سازی}

در این بخش نتیجهٔ نرمال‌سازی را برای ترم‌های برهان درست و نوع‌دار بیان
می‌کنیم: هر چنین ترمی از راه یک ترتیب کاهش مناسب به صورت بهنجار
می‌رسد. این ویژگی \emph{خاصیت نرمال‌سازی} نام دارد. از تناظر
گزاره‌ها به‌مثابهٔ نوع‌ها استفاده می‌کنیم: نرمال‌شدن ترم‌های برهان،
نرمال‌شدن !!{derivation}s استنتاج طبیعی را نتیجه می‌دهد.

نخست تابع‌هایی را تعریف می‌کنیم که پیچیدگی ترم‌ها را می‌سنجند.
اگر~$!A$ !!a{formula} باشد، \emph{طول}~$\len{!A}$ آن چنین تعریف می‌شود:
\begin{align*}
  \len{p} &= 0\\
  \len{\lfalse} &= 0\\
  \len{!A \land !B} &= \len{!A} + \len{!B} + 1\\
  \len{!A \lor !B} &= \len{!A} + \len{!B} + 1 \\
  \len{!A \lif !B} &= \len{!A} + \len{!B} + 1.
\end{align*}
پیچیدگی یک ردکس~$M$ با \emph{رتبهٔ برش}~$\cutrank{M}$ آن سنجیده
می‌شود:
\begin{align*}
  \cutrank{(\lambd[\typeof{x}{!A}][\typeof{N}{!B}])Q} &=
  \len{!A} + \len{!B} + 1 \\
  \cutrank{\ande{i}{\andi{\typeof{M}{!A}}{\typeof{N}{!B}}}} &=
  \len{!A} + \len{!B} + 1 \\
  \cutrank{\ore{\ori{i}{!A_{i}}{\typeof{M}{!A_i}}}
    {\typeof{x_1}{!A_1}}{\typeof{N_1}{!C}}
    {\typeof{x_2}{!A_2}}{\typeof{N_2}{!C}}} &=
  \len{!A_1} + \len{!A_2} + 1.
\end{align*}
پیچیدگی یک ترم برهان با پیچیده‌ترین ردکس در آن سنجیده می‌شود و اگر
ترم بهنجار باشد برابر~$0$ است:
\[
  \maxrank{M} = \max \bigl(\{\cutrank{N} \mid N
  \text{ زیرترم ردکسِ } M \text{ است}\} \cup \{0\}\bigr).
\]

\begin{lem}\ollabel{lem:subst}
  اگر $\Subst{M}{\typeof{N}{!A}}{\typeof{x}{!A}}$ یک ردکس و
  $M \not\ident x$ باشد، یکی از حالت‌های زیر برقرار است:
  \begin{enumerate}
  \item خود $M$ یک ردکس است؛ یا
  \item $M$ به شکل $\ande{i}{x}$ و $N$ به شکل
    $\andi{P_1}{P_2}$ است؛
  \item $M$ به شکل $\ore{x}{x_1}{P_1}{x_2}{P_2}$ و $N$ به شکل
    $\ori{i}{}{Q}$ است؛
  \item $M$ به شکل $x Q$ و $N$ به شکل $\lambd[y][P]$ است.
  \end{enumerate}
  در حالت نخست
  $\cutrank{\Subst{M}{N}{x}} = \cutrank{M}$؛ در حالت‌های دیگر
  $\cutrank{\Subst{M}{N}{x}} = \len{!A}$.
\end{lem}
\begin{proof}
  بر~$M$ استقرا می‌کنیم.
  \begin{enumerate}
  \item اگر $M$ تنها متغیر~$y$ و $y \not\ident x$ باشد،
    $\Subst{y}{N}{x}$ همان~$y$ است و بنابراین ردکس نیست.
  \item اگر $M$ به شکل $\andi{N_1}{N_2}$، یا $\lambd[x][N]$، یا
    $\ori{i}{!A}{N}$ باشد،
    $\Subst{M}{\typeof{N}{!A}}{\typeof{x}{!A}}$ نیز همان شکل را دارد
    و ازاین‌رو خودِ ترم حاصل ردکس نیست.
  \item اگر $M$ به شکل $\ande{i}{P}$ باشد، دو حالت را بررسی می‌کنیم.
    \begin{enumerate}
      \item اگر $P$ به شکل $\andi{P_1}{P_2}$ باشد، آنگاه
        $M \ident \ande{i}{\andi{P_1}{P_2}}$ ردکس است و آشکارا
        \[
        \Subst{M}{N}{x} \ident
        \ande{i}{\andi{\Subst{P_1}{N}{x}}{\Subst{P_2}{N}{x}}}
        \]
        نیز ردکس است. رتبه‌های برش برابرند.
      \item اگر $P$ تنها یک متغیر باشد، برای ردکس‌شدن جایگذاری باید
        همان~$x$ باشد و $N$ باید به شکل $\andi{P_1}{P_2}$ باشد. اکنون
        \[
        \Subst{M}{N}{x} \ident \Subst{\ande{i}{x}}{\andi{P_1}{P_2}}{x}
        \]
        را در نظر بگیرید که برابر $\ande{i}{\andi{P_1}{P_2}}$ است.
        رتبهٔ برش آن برابر $\len{!A}$ است.
    \end{enumerate}
  \end{enumerate}
  حالت‌های $\ore{P}{x_1}{N_1}{x_2}{N_2}$ و $P Q$ مشابه‌اند.
\end{proof}

\begin{lem}
  فرض کنید $M$ یک ردکس باشد، $M \redone M'$ با انقباض ردکسِ ریشه
  حاصل شود، و برای هر زیرترم ردکس سرهٔ~$N$ از~$M$ داشته باشیم
  $\cutrank{M} > \cutrank{N}$. آنگاه
  $\cutrank{M} > \maxrank{M'}$.
\end{lem}

\begin{proof}
  با جداسازی حالت‌ها.
  \begin{enumerate}
  \item اگر $M$ به شکل $\ande{i}{\andi{M_1}{M_2}}$ باشد، $M'$ همان
    $M_i$ است. چون هر زیرترم~$M_i$ زیرترم سرهٔ~$M$ نیز هست، ادعا
    برقرار است.
  \item اگر $M$ به شکل
    $(\lambd[\typeof{x}{!A}][N])\typeof{Q}{!A}$ باشد، $M'$ برابر
    $\Subst{N}{\typeof{Q}{!A}}{\typeof{x}{!A}}$ است. یک ردکس در~$M'$
    را در نظر بگیرید. یا این ردکس از ردکسی در~$N$ یا~$Q$ به ارث رسیده
    است و رتبهٔ برش متناظر آن بنا بر فرض کمتر از $\cutrank{M}$ است؛
    یا بر اثر جایگذاری پدید آمده و رتبهٔ برشش $\len{!A}$ است، که بنا
    بر تعریف از $\cutrank{(\lambd[x^{!A}][N])Q}$ کمتر است.
  \item اگر $M$ به شکل
    \[
    \ore{\ori{i}{}{\typeof{N}{!A_i}}}
        {\typeof{x_1}{!A_1}}{\typeof{N_1}{!C}}
        {\typeof{x_2}{!A_2}}{\typeof{N_2}{!C}}
    \]
    باشد، $M' \ident \Subst{N_i}{N}{\typeof{x_i}{!A_i}}$. یک ردکس
    در~$M'$ را در نظر بگیرید. یا این ردکس از ردکسی در~$N_i$ یا در ترم
    محمولهٔ~$N$ به ارث رسیده و همان رتبهٔ برش را دارد، که بنا بر فرض
    از $\cutrank{M}$ کمتر است؛ یا رتبهٔ برش آن $\len{!A_i}$ است، که
    بنا بر تعریف از رتبهٔ برش ردکس نمایشی بالا کمتر است.
  \end{enumerate}
\end{proof}

\begin{thm}
  هر ترم برهان درست و نوع‌دار یک دنبالهٔ کاهش متناهی به صورت بهنجار
  دارد؛ و هر !!{derivation}s نوع‌دار را می‌توان به !!{derivation}s
  بهنجار کاهش داد.
\end{thm}

\begin{proof}
  استقرای رتبهٔ بیشینه در متن منبع، اثبات معتبری برای این قضیه نیست:
  \begin{enumerate}
  \item آن استقرا پیشنهاد می‌کند ردکسی با رتبهٔ بیشینه و بدون زیرردکس
    هم‌رتبه انتخاب و تعداد ردکس‌های دارای آن رتبه شمرده شود.
  \item اما این شمارش به دو دلیل تضمین‌شده کاهش نمی‌یابد:
    \begin{enumerate}
    \item انقباض می‌تواند ردکس تازه‌ای بسازد که سراسر بافت پیرامون و
      ترمِ حاصل را در بر می‌گیرد و لم پیشین آن را کنترل نمی‌کند؛
    \item کاهش حذف جمع می‌تواند با جایگذاری، ردکسی با رتبه‌ای حتی
      بالاتر پدید آورد.
    \end{enumerate}
  \end{enumerate}
  بنابراین استقرای منبع را به‌عنوان اثبات بازنمی‌آوریم. اثبات درست
  معمولاً با نامزدهای کاهش‌پذیری یا استدلالی هم‌ارز انجام می‌شود، اما
  چنین اثباتی در اینجا ارائه نشده است؛ پس قضیه در این بخش بدون اثبات
  باقی می‌ماند.
\end{proof}

نتیجهٔ نرمال‌سازی ضعیف، پس از اثبات درست، می‌گوید هر ترم نوع‌دار
\emph{دست‌کم یک} دنبالهٔ کاهش پایان‌پذیر و یک صورت بهنجار دارد. اگر
کاهش ترم‌های~$\lambd$ را اجرای محاسبه و صورت بهنجار را مقدار محاسبه
بدانیم، این نتیجه تنها وجود یک راهبرد محاسبهٔ پایان‌پذیر را تضمین
می‌کند، نه پایان‌یافتن هر ترتیب دلخواه از کاهش‌ها. حفظ نوع نیز تضمین
می‌کند مقدار نوع درست داشته باشد. برای نمونه، اگر $N$ نوع
$!A \lif !B$ داشته باشد و آن را بر ترم~$M$ از نوع~$!A$ اعمال کنیم،
راهبرد نرمال‌ساز $(NM)$ را به ترمی چون~$N'$ از نوع~$!B$ کاهش می‌دهد.

ادعای قوی‌تر، \emph{نرمال‌سازی قوی}، قضیه‌ای جداگانه است: هیچ دنبالهٔ
نامتناهی
$N \redone N_1 \redone N_2 \redone \dots$
از کاهش‌های یک‌گامی برای ترم نوع‌دار وجود ندارد. این ادعا از برهان
بالا نتیجه نمی‌شود و در اینجا نیز برای آن اثبات یا ارجاعی عرضه نشده
است. پس نباید نتیجه گرفت هر شیوهٔ انتخاب زیرترم‌ها سرانجام پایان
می‌یابد.

از کجا می‌دانیم شیوه‌های متفاوت اجرای محاسبه \emph{همان} مقدار را
تولید می‌کنند؟ حفظ نوع فقط می‌گوید همهٔ مقدارهای حاصل نوع یکسانی دارند.
خاصیت مهم دیگر حساب~$\lambd$ نوع‌دار \emph{همگرایی} یا خاصیت
\emph{چرچ--راسر} است. اگر $N \red N_1$ و $N \red N_2$، ترمی چون
~$N'$ وجود دارد که $N_1 \red N'$ و $N_2 \red N'$. به یاد آورید
$N \red N'$ یعنی $N'$ با صفر یا چند گام کاهش~$\redone$ از~$N$ حاصل
می‌شود. اگر $N_1$ بهنجار باشد، تنها ترم~$N'$ که
$N_1 \red N'$ همان $N_1$ است (با صفر گام). بنابراین اگر $N_1$ و
$N_2$ هر دو بهنجار و رابطه همگرا باشد، هر دو برابرند. در صورت داشتن
نرمال‌سازی قوی و همگرایی، هر ترتیب کاهش پایان می‌یابد و صورت بهنجار،
مستقل از شیوهٔ محاسبه، یکتا است.

اثبات همگرایی یک رابطه معمولاً دشوار است. بااین‌حال، شرط محلی ضعیف‌تر
زیر را می‌توان مستقیماً بررسی کرد:

\begin{prop}[خاصیت ضعیف چرچ--راسر]
اگر $N \redone N_1$ و $N \redone N_2$، ترمی چون $N'$ وجود دارد که
$N_1 \red N'$ و $N_2 \red N'$.
\end{prop}

\begin{proof}
  دو انقباض را که به $N_1$ و $N_2$ می‌انجامند در نظر بگیرید. اگر
  همان وقوع ردکس را منقبض کنند، $N_1=N_2$ و ادعا فوری است. اگر دو
  وقوع در جایگاه‌های جدا از هم باشند، انقباض‌ها جابه‌جا می‌شوند:
  با انجام انقباض باقی‌مانده در هر یک از $N_1$ و $N_2$ به یک ترم
  مشترک می‌رسیم.

  فقط حالت تودرتو باقی می‌ماند. برحسب قاعدهٔ ردکس بیرونی بررسی
  می‌کنیم. برای یک تصویر مانند $\proj{1}{\pair{O_1}{O_2}}$، اگر ردکس
  در مؤلفهٔ نگه‌داشته‌شده باشد و $O_1 \redone O_1'$، دو مسیر به
  $O_1'$ می‌رسند؛ اگر در مؤلفهٔ حذف‌شده باشد، هر دو مسیر به~$O_1$
  می‌رسند. حالت تصویر دوم مشابه است. در ردکس
  $(\lambd[x][P])Q$، ردکس درونی یا در~$P$ است یا در~$Q$. سازگاری
  جایگذاریِ بی‌گرفتاری با کاهش نشان می‌دهد در حالت نخست با کاهش باقیمانده
  در $\Subst{P}{Q}{x}$، و در حالت دوم با کاهش همهٔ نسخه‌های باقیماندهٔ
  محموله در جایگذاری، به همان حاصل می‌رسیم؛ نسخهٔ پاک‌شده به هیچ گامی
  نیاز ندارد. برای حذف جمع نیز همین استدلال برقرار است: باقیماندهٔ
  شاخهٔ نگه‌داشته‌شده بعداً منقبض می‌شود، شاخهٔ حذف‌شده گامی نمی‌خواهد،
  و همهٔ نسخه‌های یک ردکس تکثیرشده در جایگذاری منقبض می‌شوند. پس در
  همهٔ حالت‌های همسان، جدا و تودرتو یک حاصل مشترک وجود دارد.
\end{proof}

\begin{lem}[لم نیومن]
اگر رابطهٔ یک‌گامی~$\redone$ نرمال‌ساز قوی و دارای خاصیت ضعیف
چرچ--راسر باشد، بستار بازتابی--تعدی آن، یعنی~$\red$، همگرا است.
\end{lem}

\begin{proof}
  چون $\redone$ نرمال‌ساز قوی است، هر ترم دست‌کم یک صورت بهنجار دارد.
  فرض کنید برخلاف ادعا ترمی چون~$N$ دو صورت بهنجار متمایز~$U$ و~$V$
  داشته باشد. چون $U\neq V$، هیچ‌یک از دو مسیر نمی‌تواند هم‌زمان مسیر
  صفرگامی مشترک باشد؛ نخستین گام‌ها را جدا می‌کنیم:
  \begin{align*}
  N &\redone N_1 \red U,\\
  N &\redone N_2 \red V.
  \end{align*}
  خاصیت ضعیف چرچ--راسر ترمی چون~$P$ می‌دهد که
  $N_1 \red P$ و $N_2 \red P$. به‌سبب نرمال‌سازی قوی، $P$ به یک
  صورت بهنجار~$W$ کاهش می‌یابد. در نتیجه $N_1$ صورت‌های بهنجار
  $U$ و~$W$ را دارد و $N_2$ صورت‌های بهنجار $V$ و~$W$ را. چون
  $U\neq V$، دست‌کم یکی از $U\neq W$ یا $V\neq W$ برقرار است. پس از
  ترمی با دو صورت بهنجار متمایز، در یک گام به ترم دیگری با دو صورت
  بهنجار متمایز رسیده‌ایم. تکرار این استدلال دنباله‌ای نامتناهی از
  گام‌های~$\redone$ می‌سازد، که با نرمال‌سازی قوی ناسازگار است.

  بنابراین صورت بهنجار هر ترم یکتا است. اگر
  $N \red N'$ و $N \red N''$، نرمال‌سازی قوی $N'$ و $N''$ را به
  صورت‌های بهنجار می‌رساند؛ یکتایی نشان می‌دهد این صورت‌ها یکی‌اند و
  در نتیجه $N'$ و $N''$ به یک ترم مشترک کاهش می‌یابند. پس~$\red$
  همگرا است.
\end{proof}


\end{document}
